% Generated draft: requires mathematical and source review.
\documentclass[11pt]{article}
\usepackage[margin=25mm]{geometry}
\usepackage{fontspec}
\setmainfont{Latin Modern Roman}
\usepackage{amsmath,amssymb,mathtools,unicode-math}
\setmathfont{Latin Modern Math}
\usepackage{enumitem,needspace,xurl,hyperref}
\setlength{\parindent}{0pt}
\setlength{\parskip}{5pt}
\setlength{\emergencystretch}{3em}
\newcommand{\N}{\mathbb{N}}
\newcommand{\Z}{\mathbb{Z}}
\newcommand{\Q}{\mathbb{Q}}
\newcommand{\R}{\mathbb{R}}
\newcommand{\C}{\mathbb{C}}
\newcommand{\F}{\mathbb{F}}
\newcommand{\A}{\mathbb{A}}
\newcommand{\PP}{\mathbb{P}}
\newcommand{\E}{\mathbb{E}}
\newcommand{\eps}{\varepsilon}
\newcommand{\dd}{\,\mathrm{d}}
\begin{document}
\section*{UnsolvedMath ID: 30003109}
\textbf{Source problem code:} catalogue entry
\subsection{Definitions and mathematical statement}
Let \(K\) be a number field and let \(g\geq 1\) be an integer. For an abelian variety \(A/K\), its Mordell--Weil rank is
\[
r_K(A)=\operatorname{rank}_{\Z}\bigl(A(K)/A(K)_{\mathrm{tors}}\bigr).
\]
An abelian variety is simple over \(K\) if it has no nonzero proper abelian subvariety defined over \(K\); geometric simplicity means simplicity after extension to an algebraic closure. The supplied record does not specify which meaning of simplicity is intended.

To explain the regulator qualification, if a symmetric ample line bundle \(L\) on \(A\), defined over \(K\), is chosen, its N\'eron--Tate height gives the pairing
\[
\langle P,Q\rangle_L=\frac{1}{2}\bigl(\widehat h_L(P+Q)-\widehat h_L(P)-\widehat h_L(Q)\bigr).
\]
For representatives \(P_1,\ldots,P_r\) of a \(\Z\)-basis of \(A(K)/A(K)_{\mathrm{tors}}\), where \(r=r_K(A)\geq1\), the corresponding regulator is
\[
\operatorname{Reg}_L(A/K)=\det\bigl(\langle P_i,P_j\rangle_L\bigr)_{1\leq i,j\leq r}.
\]
This determinant is independent of the chosen free basis, but depends on the height data and normalization. The supplied record does not specify a polarization, a line bundle, a rule for choosing one, or a normalization. This definition therefore explains the missing data and does not identify the intended source regulator.

\textbf{Conditional finiteness statement recorded by the source.} Assume the Lang--Silverman height conjecture in the form required by the cited proposition. For every number field \(K\), every integer \(g\geq1\), every integer rank bound \(M\geq1\), and every real regulator bound \(B>0\), the intended assertion is that
\[
\#\left(
\left\{A/K:\ \dim A=g,\ A\text{ simple},\ 1\leq r_K(A)\leq M,\ \operatorname{Reg}(A/K)\leq B\right\}\big/\!\sim_{\Q}
\right)<\infty.
\]
Here \(\operatorname{Reg}\), the meaning of simplicity, and \(\sim_{\Q}\) retain the unresolved source conventions. In particular, \(\sim_{\Q}\) records the supplied phrase ``\(\Q\)-isomorphism classes''; it has not been replaced by \(K\)-isomorphism. For varieties specified only over \(K\), that phrase needs an explanation of the objects and morphisms being compared. The exact Lang--Silverman hypothesis, including its uniformity and height conventions, is also absent. Thus the display records the intended quantified form, but these missing conventions prevent a fully determined mathematical proposition.

The cleaned statement uses bounded positive rank. The original extraction instead says both that the rank \(m\geq1\) is fixed and that the rank is bounded by \(m\). Finiteness for every fixed positive rank and finiteness under every finite positive rank bound are equivalent by a finite union, provided all other conventions and bounds agree.

The supplied dated assessment describes the displayed conditional implication as proved and its Lang--Silverman hypothesis as unresolved. It labels the corresponding unconditional extracted problem open. The explicit statements supplied here retain the conjectural hypothesis; removing it would be a separate unconditional target, not an established conclusion of this draft.
\subsection{Short English statement}
The record asserts conditional finiteness for simple abelian varieties of fixed dimension over a fixed number field, with positive Mordell--Weil rank bounded above and regulator bounded above, assuming Lang--Silverman. Its regulator conventions, simplicity convention, exact conjectural hypothesis, and meaning of \(\Q\)-isomorphism require clarification. The supplied curation calls the conditional implication proved and an unconditional extracted problem open; neither assessment is independently verified here.
\subsection{Sources}
\begin{itemize}
\item Supplied local UnsolvedMath registry record, permanent ID \(30003109\), title ``Regulator Finiteness from Lang--Silverman Heights,'' Oberwolfach Reports Open Problems collection, exact source problem code \texttt{catalogue entry}. No individual proposer is supplied; the recorded proposed year is 2016.
\item Supplied source citation: \emph{Diophantische Approximationen} (2017), Oberwolfach Reports, pp.~1099--1187, DOI \texttt{10.4171/owr/2016/21}. \url{https://doi.org/10.4171/owr/2016/21}. The record describes Oberwolfach Reports as documenting mathematical workshops and their open-problem discussions.
\item Supplied literature source, labelled ``OWR report'': \url{https://ems.press/content/serial-article-files/46625?nt=1}.
\item Supplied literature source, labelled ``Silverman notes'': \url{https://swc-math.github.io/aws/2024/2024SilvermanNotes.pdf}.
\item Literature assessment date, copied exactly: \texttt{2026-08-21}. This is supplied dated curation, not independently verified. It labels the unconditional extracted problem open, describes the displayed conditional implication as proved, and describes its Lang--Silverman hypothesis as unresolved. The supplied background explicitly requests independent verification of time-sensitive status and cited claims.
\item The registry assigns the statement status \texttt{corrected\_verified} and reports that the official report was checked and the rank wording corrected. These are supplied verification claims, not verification performed for this draft.
\end{itemize}
\end{document}
